\subsection{The Immense Value of Nonbuyer Interviews}

Recent buyers confirm what your marketing already does well. Nonbuyers: people who experienced the acute struggle, thoroughly evaluated your category, and explicitly chose to walk away: define the \textbf{true boundary of your market}.

Inquire without defensiveness:
\begin{itemize}[leftmargin=1.2em,itemsep=0.2ex]
  \item At what exact moment did our category become unviable or untrustworthy for you?
  \item What workaround did you return to, and how is it performing?
  \item What internal organizational risk would you have incurred by proceeding?
  \item What single condition would have had to change for you to commit?
\end{itemize}
Nonbuyer interviews routinely expose hidden switching barriers: unyielding security protocols, internal political conflicts, academic calendar mismatches, or simple lack of protected study hours.

\subsection{The Lost-Deal Forensic Audit}

When a deal collapses, sales teams often record generic CRM labels like ``lost on price.'' A neutral founder interview frequently uncovers the real cause: the engineering lead feared the curriculum lacked technical depth, the procurement office refused a 100\% upfront advance, or the sponsor lost their internal political champion.

Compare the prospect's post-mortem account against contemporaneous email threads, proposal revisions, and evaluation scorecards. Distinguish between preventable sales mistakes and \textbf{fundamental structural disqualifications}. If a customer walked away because they demanded an accredited university degree, no amount of product polish or price discounting could have saved the transaction.

\subsection{A Bad Interview vs. A Disciplined Interview}

\noindent\textbf{The Naive, Biased Interview:}
\begin{quote}
\small
\textbf{Founder:} ``We're building an advanced, hands-on, live studio for engineers to master production AI systems. Do you think that's needed today?''\\
\textbf{Prospect:} ``Oh yes, definitely. AI is everywhere and everyone needs upskilling.''\\
\textbf{Founder:} ``Awesome. We're thinking of charging ₹45,000 for 8 weeks of intensive code reviews. Does that sound reasonable for a career accelerator?''\\
\textbf{Prospect:} ``Sure, for people who want to level up, that sounds very reasonable.''\\
\textbf{Founder:} ``Fantastic! Can I add you to our priority waitlist?''\\
\textbf{Prospect:} ``Yes, absolutely, send me an email when you launch.''
\end{quote}
\textbf{Audit:} Zero empirical information gained. The founder asked leading questions, fished for compliments, proposed hypothetical pricing, and secured a costless verbal promise.

\vspace{0.8ex}
\noindent\textbf{The Forensic, Behavioral Interview:}
\begin{quote}
\small
\textbf{Founder:} ``When was the last time you paid money or took personal time to learn a new software framework for work?''\\
\textbf{Prospect:} ``About 9 months ago. I bought an online course on Kubernetes.''\\
\textbf{Founder:} ``What specifically happened at work that made you buy it?''\\
\textbf{Prospect:} ``Our DevOps lead quit unexpectedly, and my manager asked me to take over the deployment clusters.''\\
\textbf{Founder:} ``Did you finish the course?''\\
\textbf{Prospect:} ``I watched the first 3 modules, then hit an obscure network configuration bug. The video didn't help. I spent 4 nights debugging it with a colleague on Slack.''\\
\textbf{Founder:} ``Did your employer reimburse the course fee?''\\
\textbf{Prospect:} ``No, our company requires a formal vendor agreement and manager approval 30 days in advance, so I just paid the ₹3,000 out of my own pocket to save time.''
\end{quote}
\textbf{Audit:} High empirical yield. The founder learned the trigger (colleague departure), the failure mode of video courses (lack of debugging feedback on edge cases), the true reliance on peer assistance, and the friction of corporate reimbursement rules.

\subsection{What Never to Do in an Interview}

Strictly avoid the following fatal practices:
\begin{itemize}[leftmargin=1.2em,itemsep=0.2ex]
  \item Never pitch your solution before thoroughly understanding their past behavior.
  \item Never ask: ``Would you buy a product that did X?''
  \item Never argue with a participant or defend your product concept.
  \item Never ask leading questions: ``Don't you agree that existing courses are superficial?''
  \item Never record without explicit, recorded consent.
  \item Never harvest confidential company code, private customer data, or proprietary trade secrets.
\end{itemize}

\subsection{Contextual Observation: Actions Over Narratives}

Human memory is reconstructive and prone to rationalization. When interviewees narrate their work, they omit mundane workarounds, subconscious habits, and embarrassing mistakes. \textbf{Contextual inquiry}: observing the user perform the actual task in their live working environment: reveals the unvarnished reality [R11].

Watch an engineer evaluate an AI model's output:
\begin{itemize}[leftmargin=1.2em,itemsep=0.2ex]
  \item Notice which tools they open first and which terminal commands fail.
  \item Track how many times they copy-paste text between browser tabs and internal editors.
  \item Observe how they verify whether an answer is correct: do they inspect raw tokens, query a database, or rely on visual guesswork?
\end{itemize}
These micro-frictions expose the exact points where human capability and software tools break down.

\subsection{The Artifact Method: Analyzing Behavioral Traces}

Physical and digital artifacts provide objective, corroborating evidence of customer struggle. Ask prospects to share (under appropriate non-disclosure and redaction protocols):
\begin{itemize}[leftmargin=1.2em,itemsep=0.2ex]
  \item Anonymized spreadsheets used to track model evaluations.
  \item Internal architecture review rejection memos.
  \item CI/CD pipeline error logs showing test failure rates.
  \item Procurement guidelines detailing mandatory vendor requirements.
\end{itemize}
An artifact cannot lie about what occurred. It grounds discovery in observable historical artifacts.

\subsection{Diary Studies and Longitudinal Observation}

For intermittent, sporadic problems, a single retrospective interview can miss cyclical patterns. In a \textbf{diary study}, a small cohort of qualified practitioners logs events over 2 to 4 weeks:
\begin{itemize}[leftmargin=1.2em,itemsep=0.2ex]
  \item Exact time and trigger when an AI evaluation defect occurred.
  \item Tools deployed and external resources consulted.
  \item Total minutes spent attempting resolution before giving up or asking peers.
  \item Emotional frustration and operational delay incurred.
\end{itemize}
Keep logging frictionless (e.g., a simple mobile form or private messaging channel). Diary studies expose how frequently the acute pain genuinely recurs in normal working life.

\subsection{Secondary Public Evidence: Search, Forums, and Repositories}

Public technical forums, GitHub issues, Discord developer communities, Stack Overflow questions, and Reddit discussions offer accessible preliminary clues. They reveal the specific technical terminology practitioners use when struggling with a new framework.

However, recognize the severe selection bias of public forums: the loudest voices represent extreme enthusiasts or frustrated beginners; satisfied, competent corporate engineers rarely post public complaints. Use public discussions to generate hypotheses, never to quantify market demand.

\subsection{Customer Support as Continuous Field Discovery}

Once a product is live, customer support tickets, refund requests, and withdrawal emails represent the most honest, unfiltered customer research available.

A refund request stating ``I didn't have enough time'' is not an excuse to dismiss: it is diagnostic evidence of an unresolved onboarding bottleneck, an inaccurate prerequisite filter, or an overwhelming weekly workload. Categorize every failure mode in a live \textbf{Support Defect Ledger}. Feed these findings directly into admissions criteria and curriculum architecture.
